\documentclass[11pt, a4paper]{article}

\usepackage[T1]{fontenc}
\usepackage[utf8]{inputenc}
\usepackage[a4paper, top=2.5cm, bottom=2.5cm, left=2.5cm, right=2.5cm]{geometry}
\usepackage{booktabs}
\usepackage{enumitem}
\usepackage[hidelinks]{hyperref}
\usepackage{parskip}
\usepackage{amsmath}
\usepackage{listings}
\usepackage{xcolor}

\definecolor{codigobg}{RGB}{245,245,245}
\definecolor{keyword}{RGB}{0,100,150}
\lstset{language=R, basicstyle=\ttfamily\small, backgroundcolor=\color{codigobg},
  keywordstyle=\color{keyword}\bfseries, frame=single, breaklines=true, showstringspaces=false}

\setlength{\parindent}{0pt}

\pagestyle{plain}

\begin{document}

\begin{center}
{\large\bfseries PONTIFICIA UNIVERSIDAD JAVERIANA}\\[3pt]
Facultad de Ciencias Econ\'omicas y Administrativas\\[2pt]
Estad\'istica 2026-II --- Prof.\ Jaime Polanco Jim\'enez\\[6pt]
{\LARGE\bfseries Problem Set 4}\\[3pt]
{\large Intervalos de Confianza}\\[6pt]
Asignado: Septiembre 11 (Sesi\'on 4) \quad$|$\quad Entrega: Septiembre 18 (Sesi\'on 5)
\end{center}

\vspace{6pt}
\rule{\linewidth}{0.5pt}
\vspace{2pt}
{\centering\small Repositorio: \url{https://github.com/polanco-jaime/Curso-Estadistica-2026-3}\par}
\vspace{6pt}

\textbf{Instrucciones.} Este Problem Set es individual y se resuelve en R sobre
microdatos ICFES. La entrega es digital al inicio de la clase indicada.
El c\'odigo R debe incluirse.

\vspace{10pt}

\section*{Enunciado}

Utilizando la variable \texttt{MOD\_INGLES\_PUNT} (puntaje de ingl\'es en Saber Pro)
para estudiantes de Negocios Internacionales (NI), construya intervalos de confianza
y analice la precisi\'on de la estimaci\'on.

\begin{enumerate}[label=\alph*)]
  \item \textbf{IC para la media global:} calcule un intervalo de confianza al
        95\% para la media poblacional de \texttt{MOD\_INGLES\_PUNT} usando la
        distribuci\'on $t$ de Student. Use \texttt{t.test()}.
        \begin{itemize}
          \item Reporte el intervalo $[\,\underline{x},\, \overline{x}\,]$.
          \item Interprete el resultado en contexto.
          \item ¿Cu\'al es el margen de error?
        \end{itemize}

  \item \textbf{IC por departamento:} calcule intervalos de confianza al 95\%
        para la media de \texttt{MOD\_INGLES\_PUNT} en cada uno de los 33 departamentos
        de Colombia (variable \texttt{ESTU\_DEPTO\_RESIDE}).
        \begin{itemize}
          \item Presente una tabla con: departamento, media, l\'imite inferior,
                l\'imite superior y tama\~no de muestra $n$.
          \item ¿Qu\'e departamento tiene el IC m\'as amplio? ¿Por qu\'e?
          \item ¿Qu\'e departamento tiene el IC m\'as estrecho? ¿Por qu\'e?
        \end{itemize}

  \item \textbf{Forest plot:} grafique un \textit{forest plot} con los 33 intervalos
        de confianza departamentales. El gr\'afico debe incluir:
        \begin{itemize}
          \item Puntos para las medias de cada departamento.
          \item Barras horizontales para los intervalos de confianza.
          \item Una l\'inea vertical de referencia en la media nacional.
          \item Etiquetas de departamento en el eje $y$.
        \end{itemize}
        ¿Qu\'e departamentos se distinguen significativamente de la media nacional?

  \item \textbf{IC para una proporci\'on:} calcule un intervalo de confianza al
        95\% para la proporci\'on de estudiantes de NI que superan el nivel B1 en
        ingl\'es (es decir, \texttt{MOD\_INGLES\_DESEM} $\in \{\text{B1, B+}\}$).
        Use \texttt{prop.test()}.
        \begin{itemize}
          \item Reporte el intervalo $[\,\underline{p},\, \overline{p}\,]$.
          \item Interprete el resultado. ¿Qu\'e porcentaje de estudiantes NI alcanza
                B1 o superior?
        \end{itemize}

  \item \textbf{Tama\~no de muestra \'optimo:} calcule el tama\~no de muestra $n$
        necesario para estimar la media de \texttt{MOD\_INGLES\_PUNT} con un margen
        de error de $\pm 2$ puntos al 95\% de confianza. Use la f\'ormula:
        \[
          n = \left( \frac{z_{\alpha/2} \cdot \sigma}{E} \right)^2
        \]
        donde $z_{0.025} = 1.96$ y $\sigma$ es la desviaci\'on est\'andar muestral.
        ¿Cu\'antos estudiantes se necesitar\'ian?

  \item \textbf{Conclusiones:} resuma en un p\'arrafo los principales hallazgos
        sobre la precisi\'on de las estimaciones y las diferencias departamentales.
\end{enumerate}

\vspace{10pt}
\rule{\linewidth}{0.4pt}
\begin{center}
{\small Cada entrega completa suma $+0.0625$ a la nota final.}
\end{center}

\end{document}
